% =====================================================================
%  FICHE 8 — THÈME 4 — NIVEAU DIFFICILE — CORRIGÉ
% =====================================================================
\documentclass[11pt,a4paper]{article}
\usepackage[utf8]{inputenc}
\usepackage[T1]{fontenc}
\usepackage{lmodern}
\usepackage[french]{babel}
\usepackage{amsmath,amssymb,mathtools}
\usepackage{icomma}
\usepackage{siunitx}
\sisetup{output-decimal-marker={,},per-mode=power,group-separator={\,},group-minimum-digits=5}
\usepackage[version=4]{mhchem}
\usepackage{xcolor}
\usepackage{tikz}
\usepackage[most]{tcolorbox}
\usepackage{enumitem}
\usepackage{array}
\usepackage{fancyhdr}
\usepackage[hidelinks]{hyperref}
\usetikzlibrary{arrows.meta,positioning,calc}
\usepackage[a4paper,margin=2.1cm,headheight=26pt,headsep=10pt]{geometry}
\usepackage{titlesec}

\definecolor{primary}{RGB}{150,98,162}
\definecolor{primarydark}{RGB}{92,52,110}
\definecolor{excol}{RGB}{0,135,90}

\tcbset{boxbase/.style={enhanced,breakable,boxrule=0.9pt,arc=2.5pt,left=3mm,right=3mm,top=2.3mm,bottom=2.3mm,fonttitle=\bfseries,coltitle=white,toptitle=0.6mm,bottomtitle=0.6mm}}
\newtcolorbox[auto counter]{correction}[1][]{boxbase,colback=excol!4,colframe=excol,title={\textsc{Corrigé}~\thetcbcounter\ \ifx\relax#1\relax\relax\else\textmd{--- \itshape #1}\fi}}

\newcommand{\dS}{\Delta S}
\newcommand{\dH}{\Delta H}
\newcommand{\rep}[1]{\textcolor{primarydark}{\textbf{#1}}}

\pagestyle{fancy}
\fancyhf{}
\fancyhead[L]{\small\textcolor{primarydark}{\textbf{Fiche 8} — Thème 4 : Entropie et deuxième principe}}
\fancyhead[R]{\small\textcolor{primarydark}{\textsc{Difficile — corrigé}}}
\fancyfoot[C]{\small\thepage}
\renewcommand{\headrulewidth}{0.8pt}
\renewcommand{\headrule}{\hbox to\headwidth{\color{primary}\leaders\hrule height \headrulewidth\hfill}}
\setlength{\parindent}{0pt}\setlength{\parskip}{3pt}

\begin{document}

\begin{center}
\begin{tikzpicture}
\node[rectangle,rounded corners=7pt,fill=excol,text=white,minimum width=16cm,
      minimum height=1.1cm,align=center,inner sep=6pt]
  {{\large\bfseries Thème 4 — Entropie et deuxième principe}\\[1pt]{\small Niveau \textsc{Difficile} — corrigé détaillé}};
\end{tikzpicture}
\end{center}

\begin{correction}[comparer l'ampleur du désordre]
\begin{enumerate}[label=\alph*),leftmargin=2em,topsep=2pt,itemsep=2pt]
  \item (i) apparition de gaz : $\dS>0$ ; (ii) $2\to1$ mol de gaz : $\dS<0$ ;
        (iii) solide $\to$ gaz : $\dS>0$ ; (iv) gaz $\to$ solide : $\dS<0$ ;
        (v) dissolution : $\dS>0$.
  \item La plus grande \emph{augmentation} est \rep{(iii)} : passer d'un solide très ordonné directement
        à un gaz est le saut de désordre le plus important (bien plus qu'ajouter une mole de gaz ou
        disperser des ions).
  \item La plus grande \emph{diminution} est \rep{(iv)} : condenser un gaz en solide ordonné détruit le
        maximum de désordre (davantage que perdre une mole de gaz en (ii)).
\end{enumerate}
\end{correction}

\begin{correction}[la poche de froid : endothermique mais spontanée]
\begin{enumerate}[label=\alph*),leftmargin=2em,topsep=2pt,itemsep=2pt]
  \item La poche \emph{refroidit} : la dissolution \textbf{absorbe} de la chaleur au milieu, elle est
        \rep{endothermique}, donc \rep{$\dH>0$}.
  \item Un solide ordonné se disperse en ions dans l'eau : le désordre augmente, \rep{$\dS_{\text{système}}>0$}.
  \item Le sens spontané est celui pour lequel $\dS_{\text{univers}}>0$. Ici, la forte augmentation de
        désordre du système ($\dS_{\text{système}}>0$) l'emporte : \rep{$\dS_{\text{univers}}>0$}, la
        dissolution est donc spontanée \emph{bien qu'elle absorbe de l'énergie}.
  \item En ne regardant que l'énergie, on prédirait \emph{à tort} que la réaction ne se produit pas
        (puisqu'elle « coûte » de la chaleur). \rep{Le sens d'évolution n'est pas dicté par la seule
        énergie} : le désordre (l'entropie) est un moteur à part entière.
\end{enumerate}
\end{correction}

\begin{correction}[sublimation de la neige carbonique et principes]
\begin{enumerate}[label=\alph*),leftmargin=2em,topsep=2pt,itemsep=2pt]
  \item Solide $\to$ gaz : \rep{$\dS_{\text{système}}>0$} (fort).
  \item La sublimation est spontanée à température ambiante, donc \rep{$\dS_{\text{univers}}>0$}.
  \item Le $\ce{CO2}$ \emph{solide} est un réseau cristallin ordonné (molécules quasi immobiles) ; le
        $\ce{CO2}$ \emph{gazeux} est très désordonné (molécules libres, dispersées) : son entropie est
        \rep{bien plus grande}.
  \item Troisième principe : à $T=\SI{0}{\kelvin}$, l'entropie du cristal parfait est
        \rep{$S=\SI{0}{\joule\per\mol\per\kelvin}$}. Ce « zéro » commun donne aux entropies une
        \rep{valeur absolue}, alors que pour l'enthalpie on ne mesure que des \emph{variations} $\dH$
        (pas de $H$ absolu).
\end{enumerate}
\end{correction}

\end{document}
